\section{2021 Geometry \& Topology}

\subsection{Exam}

\begin{exercise}
Let $S^{n}$ be the unit sphere in $\mathbb{R}^{n+1}$.
(a) Find a 6-form $\alpha$ on $\mathbb{R}^{7} \setminus \{0\}$ such that $d\alpha = 0$, and $\int_{S^{6}} \alpha = 1$.
(b) For any smooth map $f : S^{11} \to S^{6}$, show that $f^{*}\alpha = d\varphi$.
(c) Let $H(f) = \int_{S^{11}} \varphi \wedge d\varphi$. Show that $H(f)$ is independent of $\varphi$.
(d) Show that $H(f)$ is an even integer.
\end{exercise}

\begin{solution}
\textbf{Pedagogical Motivation:} This problem constructs the \textit{Hopf invariant}. For a map $f: S^{2n-1} \to S^n$, the Hopf invariant measures how much the preimages of two points in $S^n$ are "linked" in $S^{2n-1}$. 

\textbf{Formal Proof:}
(a) Let $\omega$ be the volume form of $S^6$. The form $\alpha = \frac{1}{\text{Vol}(S^6)} \frac{\sum (-1)^{i-1} x^i dx^1 \wedge \dots \wedge \widehat{dx^i} \wedge \dots \wedge dx^7}{r^7}$ is closed and integrates to 1.
(b) Since $H^6(S^{11}) = 0$, the pullback $f^* \alpha$ (which is closed) must be exact. Thus $f^* \alpha = d\varphi$.
(c) If $\varphi'$ is another choice, $\varphi' - \varphi$ is closed, hence exact ($\varphi' - \varphi = d\eta$) because $H^5(S^{11})=0$.
\[ \int \varphi' \wedge d\varphi' - \int \varphi \wedge d\varphi = \int d\eta \wedge f^* \alpha = \int d(\eta \wedge f^* \alpha) = 0. \]
(d) \textbf{Geometric Intuition:} For $f: S^{2n-1} \to S^n$, the Hopf invariant is an integer. For $n$ even, it can be any integer if $n \in \{2, 4, 8\}$. For $n=6$, a deep result by Adams (Hopf Invariant One) says that no map has Hopf invariant 1. For $n$ even and $n \notin \{2, 4, 8\}$, the Hopf invariant must be even because the Whitehead product $[\iota_n, \iota_n]$ has Hopf invariant 2 and spans the image of the Hopf homomorphism.
\end{solution}

\begin{exercise}
Either construct a positive smooth function $h_{1}$ such that $\mathrm{Ric}(h_{1}) = 1$, or show that no such function $h_{1}$ exists on $\mathbb{R}^2$.
\end{exercise}

\begin{solution}
\textbf{Geometric Intuition:} In dimension 2, the Ricci curvature is essentially the Gaussian curvature. The problem asks for a metric $g = h dx^2 + h dy^2$ with constant Gaussian curvature $K=1$ on the whole plane. We know the sphere has $K=1$ but it is compact. The plane is non-compact and "too small" to fit a sphere's worth of curvature globally without blowing up.

\textbf{Formal Proof:}
The condition is $\Delta \log h = h$. Let $u = \log h$. We have $\Delta u = e^u$.
By the \textit{Maximum Principle} or the \textit{Ahlfors Lemma}, such a solution cannot exist on $\mathbb{R}^2$.
If it existed, its spherical average $v(r)$ would satisfy $v'' + \frac{1}{r} v' \geq e^v$. 
Comparison with the explicit solution $w(r) = \log \frac{8a^2}{(1-a^2 r^2)^2}$ (which blows up at $r=1/a$) shows that $v(r)$ must blow up at some finite radius. Thus no global solution exists on $\mathbb{R}^2$.
\end{solution}

\begin{exercise}
(a) Which one of $t\frac{\partial}{\partial x^{\alpha}}\big|_{\gamma(t)}$ or $\frac{\partial}{\partial x^{\alpha}}\big|_{\gamma(t)}$ is a Jacobi field? 
(b) Compute $\frac{\partial^{2} g_{22}}{\partial x^{1} \partial x^{1}}$ at $p$.
(c) Show $\max | \partial_1 g_{22} | \leq C \delta A$.
\end{exercise}

\begin{solution}
\textbf{Formal Proof:}
(a) $J(t) = t \frac{\partial}{\partial x^\alpha}$ is a Jacobi field. In normal coordinates, $x(t) = t e_1$. A variation of geodesics through the origin is $f(t, s) = \exp_p(t(e_1 + s e_\alpha))$. The variation vector field is $V(t) = \frac{\partial f}{\partial s}(t, 0) = d(\exp_p)_{t e_1} (t e_\alpha) = t \frac{\partial}{\partial x^\alpha}$. Any variation of geodesics yields a Jacobi field.
(b) In normal coordinates, $g_{ij} = \delta_{ij} - \frac{1}{3} R_{ikjl} x^k x^l + O(|x|^3)$.
Differentiating twice with respect to $x^1$:
\[ \frac{\partial^2 g_{22}}{\partial (x^1)^2} = -\frac{2}{3} R_{2121} = \frac{2}{3} R_{1212} = -\frac{2}{3} K(e_1, e_2). \]
(c) Since $g_{22} = 1 - \frac{t^2}{3} R_{1212} + O(t^3)$, the derivative is $\partial_1 g_{22} = -\frac{2t}{3} R_{1212} + \dots$. This is bounded by $C t \|R\|$.
\end{solution}

\begin{exercise}
(a) Show $SO(n)$ is compact. (b) Is $SO(3) \cong \mathbb{RP}^3$? (c) $\pi_1(SO(2020))$.
\end{exercise}

\begin{solution}
(a) $SO(n)$ is closed ($A^T A = I$ is a continuous condition) and bounded ($\sum a_{ij}^2 = n$). By Heine-Borel, it is compact.
(b) Yes. The double cover $S^3 \to SO(3)$ (using quaternions) has kernel $\{\pm 1\}$. Thus $SO(3) \cong S^3/\{\pm 1\} = \mathbb{RP}^3$.
(c) $\pi_1(SO(2)) = \mathbb{Z}$. For $n \geq 3$, $\pi_1(SO(n)) = \mathbb{Z}_2$. Thus $\pi_1(SO(2020)) = \mathbb{Z}_2$.
\end{solution}

\begin{exercise}
Suppose $X$ is a topological space and $\pi: \mathbb{R}^2 \to X$ is a covering map. Let $K \subset X$ be a compact subset and $B \subset \mathbb{R}^2$ be the closed unit ball.
(a) If the restriction $\pi|_{\mathbb{R}^2 \setminus B}: \mathbb{R}^2 \setminus B \to X \setminus K$ is a homeomorphism, show that $\pi: \mathbb{R}^2 \to X$ is a homeomorphism.
(b) If we only assume that $\mathbb{R}^2 \setminus B$ is homeomorphic to $X \setminus K$ (not necessarily via the covering map), is $X$ necessarily homeomorphic to $\mathbb{R}^2$?
\end{exercise}

\begin{solution}
\textbf{Geometric Intuition:} A covering map between manifolds has a constant "number of sheets". If the map is a homeomorphism on the "ends" of the spaces, it means there is only one sheet at infinity. For a connected space, this forces there to be only one sheet everywhere.

\textbf{Formal Proof:}
(a) A covering map $\pi: E \to X$ between connected manifolds has a constant degree $n$, where $n$ is the number of points in the fiber $\pi^{-1}(x)$ for any $x \in X$.
The condition that $\pi|_{\mathbb{R}^2 \setminus B}$ is a homeomorphism onto $X \setminus K$ implies that for any $x \in X \setminus K$, the fiber $\pi^{-1}(x) \cap (\mathbb{R}^2 \setminus B)$ contains exactly one point.
Since $B$ is compact, its image $\pi(B)$ is compact in $X$.
Since $X \setminus K$ is homeomorphic to an annulus (which is non-compact), we can choose a point $x_0 \in X \setminus (K \cup \pi(B))$.
For this $x_0$, since $x_0 \notin \pi(B)$, all preimages $\pi^{-1}(x_0)$ must lie in $\mathbb{R}^2 \setminus B$.
From the homeomorphism property at the ends, we know there is exactly one such preimage.
Thus the degree of the cover is $n = 1$. A covering map of degree 1 is a homeomorphism.

(b) \textbf{No.} Consider the open Mobius strip $M$. The universal cover of $M$ is $\mathbb{R}^2$.
Let $K$ be a compact neighborhood of the central circle in $M$. The complement $M \setminus K$ is homeomorphic to an open annulus $S^1 \times (0, 1)$.
The complement of the unit ball in $\mathbb{R}^2$ is also an open annulus.
Thus $M \setminus K \cong \mathbb{R}^2 \setminus B$, but $M$ is not homeomorphic to $\mathbb{R}^2$ (it is non-orientable and has $\pi_1(M) \cong \mathbb{Z}$).
\end{solution}

\begin{exercise}
(a) Give an example of a graph whose fundamental group is the free group of rank 2, $F_2$.
(b) Does $F_2$ contain a proper subgroup isomorphic to $F_2$?
(c) Does $F_2$ contain a proper subgroup of finite index that is isomorphic to $F_2$?
\end{exercise}

\begin{solution}
\textbf{Formal Proof:}
(a) The figure-eight graph (a wedge of two circles $S^1 \vee S^1$) has fundamental group $F_2$. More generally, any connected graph with $E$ edges and $V$ vertices has fundamental group $F_{E-V+1}$. For $E=2, V=1$, we get $F_2$.

(b) \textbf{Yes.} Let $F_2 = \langle a, b \rangle$. The subgroup $H = \langle a^2, b^2 \rangle$ is a free group. Since it is generated by two elements that do not satisfy any relation other than those inherited from $F_2$ (which has none), it is isomorphic to $F_2$. It is proper because $a \notin H$.

(c) \textbf{No.} The rank of a finite index subgroup of a free group is given by the \textit{Nielsen-Schreier formula}:
\[ \text{rank}(H) = 1 + [G:H](\text{rank}(G) - 1). \]
Let $\text{rank}(G) = 2$ and $[G:H] = d$. Then:
\[ \text{rank}(H) = 1 + d(2 - 1) = 1 + d. \]
If $H \cong F_2$, then $\text{rank}(H) = 2$.
Substituting into the formula: $2 = 1 + d \implies d = 1$.
Thus, the index must be 1, meaning $H$ is not a proper subgroup. Any proper subgroup of $F_2$ with finite index must have rank at least 3.
\end{solution}
